\documentclass[a4paper, 11pt]{article}

\usepackage[utf8]{inputenc}
\usepackage[ngerman]{babel}  % deutsche Silbentrennung
\usepackage[T1]{fontenc}
%\usepackage[round]{natbib}
\usepackage{graphicx}
\usepackage{geometry}
\usepackage{makeidx} % Package zur Indexerstellung
\usepackage{listings}
\usepackage{url}
\parindent 0pt
\usepackage{rotating}

\usepackage{graphicx} 
\usepackage{subfigure}

\geometry{outer=25mm , inner=25mm , bottom=25mm , top=25mm }	
\newcommand{\cmd}[1]{\textbackslash #1}
\newcommand{\meta}[1]{\textless#1\textgreater}
\newcommand{\val}[1]{\{\meta{#1}\}}
\newcommand{\valo}[1]{[\meta{#1}]}
\newcommand{\incsection}[1]{\include{kapitel/#1}}
\lstset{
basicstyle=\footnotesize\ttfamily, % Standardschrift
         numbers=left,               % Ort der Zeilennummern
         numberstyle=\tiny,          % Stil der Zeilennummern
         stepnumber=1,               % Abstand zwischen den Zeilennummern
         numbersep=5pt,              % Abstand der Nummern zum Text
         tabsize=2,                  % Groesse von Tabs
         extendedchars=true,         %
         breaklines=true,            % Zeilen werden Umgebrochen
         keywordstyle=\color{red},
         stringstyle=\color{white}\ttfamily, % Farbe der String
         showspaces=false,           % Leerzeichen anzeigen ?
         showtabs=false,             % Tabs anzeigen ?
         xleftmargin=17pt,
         framexleftmargin=17pt,
         framexrightmargin=5pt,
         framexbottommargin=4pt,
         %backgroundcolor=\color{lightgray},
         showstringspaces=false,      % Leerzeichen in Strings anzeigen ? 
         numberbychapter=true
}

\begin{document}
\setlength{\parindent}{0.0cm} % Erstzeilen-Einzug
\begin{titlepage}

\large
\begin{figure}
  \begin{center}
    \hbox to \hsize{%
      \begin{tabular}[m]{c}
        \includegraphics[width=4cm]{images/hochschule_zittau_goerlitz_logo.png}
      \end{tabular}
      \hfill%
      \begin{tabular}[m]{c}
        Hochschule Zittau/Görlitz\\
        Fakultät Elektrotechnik-Informatik \\
        
      \end{tabular}%
    }
  \end{center}
\end{figure}

\begin{center}
\rule{0pt}{0pt}

\begin{huge}
Belegarbeit Bildverarbeitung und Mobile Roboter
\end{huge}
\vspace{.5cm}
\begin{center}
bei\\
\vspace{.5cm}
Prof. Bischoff und Prof. Böhm \\Hochschule Zittau/Görlitz
\end{center}


\vspace{2cm}

Autonom fahrender Roboter mit Straßenschildererkennung.\\

\vspace*{2cm}

\begin{tabular}[m]{c}
Dominik Bitterlich \\Hochschule Zittau/Görlitz\\ Matr.-Nr.: 203359
\end{tabular}
\hspace*{4cm}
\vspace*{.5cm}
\begin{tabular}[m]{c}
Alexander Häse \\Hochschule Zittau/Görlitz\\ Matr.-Nr.:42670
\end{tabular}

\vspace*{.5cm}
\begin{tabular}[m]{c}
Daniel Richter \\Hochschule Zittau/Görlitz\\ Matr.-Nr.: 202901
\end{tabular}
\hspace{4cm}
\vspace*{.5cm}
\begin{tabular}[m]{c}
Tom Schumann \\Hochschule Zittau/Görlitz\\ Matr.-Nr.: 207329
\end{tabular}

\vspace{3cm}

9. Juli 2014

\end{center}
\end{titlepage} % Referenzierung des Titelblattes
\tableofcontents  % Inhaltsverzeichnis anzeigen

\newpage
\thispagestyle{empty}
\addcontentsline{toc}{section}{Selbstständigkeitserklärung}

Hiermit erkläre ich, dass ich diese Arbeit selbstständig verfasst und keine anderen als die angegebenen Quellen und Hilfsmittel benutzt habe.
Diese Arbeit wurde bisher auch keiner anderen Prüfungsbehörde vorgelegt und auch noch nicht veröffentlicht.

\vspace{5cm}

Görlitz, den 1.2.2013 %\hspace{4cm} \line(1,0){150}

\vspace{1cm}
\line(1,0){150}\\
Der Eine

\vspace{1cm}
\line(1,0){150}\\
Der Andere

%Kapitel in orderner sections als eigene dateien
\section{Einleitung}

\section{Kommunikation zwischen Raspberry Pi und Lego EV3}
Nach dem sich abzeichnete, das der Lego EV3 noch nicht mit einer Webcam kommunizieren kann, entschied sich das Team zur Auslagerung der Bildverarbeitung auf einen Raspberry Pi. Für die Kommunikation zwischen Raspberry Pi und EV3 zu bewerkstelligen, wurden drei Wege verfolgt. Zuerst sollte die beiden Geräte über Bluetooth gekoppelt werden, dies schlug durch Probleme mit der Bluez Bibliothek fehl. Durch eine Verbindung des Raspberry Pi mit dem EV3 über USB, funktionierte die Kommunikation.
\subsection{Umsetzung}
Im Server Teil des Projekts, stellt der Raspberry Pi mit dem Befehl \lstinline$ifconfig$ eine TCP/IP-Verbindung zwischen Raspberry Pi und EV3 aufgebaut. Durch den in Java implementierten ProzessBuilder werden die benötigten Bash-Befehle aus dem Java Programm heraus ausgeführt, was zusätzliche Einstellungen vor dem Starten des Servers überflüssig machte. Nach der erfolgreichen Verbindung wird auf dem Raspberry Pi ein Socket geöffnet, über den die beiden Geräte kommunizieren. Auch die Verbindungswiederaufnahme nach der Verbindung nach Absturz des EV3-Programms ist möglich, doch sollte die Hardware Verbindung der beiden Geräte bestehen bleiben.
\section{Bedienungsanleitung}
\subsection{Benötigte Hardware}
Für das Projekt wurde folgende Hardware benötigt:
\begin{itemize}
\item Lego EV3 Education mit Akku
\item Raspberry Pi Model B
\item Raspberry PI NOIR Kamera
\item SD-Karte für Raspberry Pi (hier: 16GB Class 10, der Swap wird benötigt)
\item WLan-Stick (hier: Edimax ew-7811un)
\item Raspberry Pi Akku (hier: XORO MPB 520 mit 5200mAh)
\end{itemize}
\subsection{Vorbereitung}
Folgende Schritte müssen nur einmal durch geführt werden.\\
Ausgangspunkt ist Vorbereiteter Raspberry Pi mit
\begin{enumerate}
\item 
\end{enumerate}
\subsection{Bedienung}
\begin{enumerate}
\item USB-Verbindung zwischen Raspberry Pi und EV3 herstellen
\item via 
\end{enumerate}
%content

% TODO anpassen, DIES HIER DIENT NUR ALS BEISPIEL
\addcontentsline{toc}{section}{Quellenverzeichnis}
\renewcommand{\refname}{Quellenverzeichnis}
\def\UrlFont{\bfseries}
\nocite{*}
\begin{thebibliography}{99}
	\bibitem{1} \url{https://developers.google.com/apps-script/articles/bigquery_tutorial}
	\bibitem{3} \url{https://developers.google.com/apps-script/overview}
	\bibitem{2} \url{https://developers.google.com/bigquery/docs/overview}
	\bibitem{5} \url{https://developers.google.com/bigquery/docs/dataset-gsod}
	\bibitem{4} \url{ftp://ftp.ncdc.noaa.gov/pub/data/inventories/ISH-HISTORY.TXT}
	\bibitem{6} \url{https://en.wikipedia.org/wiki/BigQuery}

\end{thebibliography}

\addcontentsline{toc}{section}{Abbildungen}
\renewcommand{\refname}{Abbildungen}
\def\UrlFont{\bfseries}

\begin{appendix}
\section*{Abbildungen}

%\begin{figure}[htbp]
%  \begin{sideways}
%    \begin{minipage}{19cm}
%		\includegraphics[width=0.90\textwidth]{./diagramme/jahresdurchschnitt}
%		\caption{Jahresdurchschnittstemperatur}
%    \end{minipage}
%  \end{sideways}
%  \centering
%\end{figure}


%\begin{sidewaysfigure}[h]
%\centering\includegraphics[width=0.80\textwidth]{./diagramme/monatJahr}
%\caption{Monatsdurchschnittstemperatur über die Jahre}
%\end{sidewaysfigure}

\end{appendix}

\newpage
 


\end{document}
